% Schematic of the no-late-interaction event L_{T,t}.
% The visual conventions are inherited from Figure 1 (patch-geometry.tex):
% red striped patch boundaries, gray outgoing-terminal patches, thin black
% finite-patch boundaries, and black successful interactions.
\begingroup
\definecolor{patchfigred}{RGB}{175,32,38}
\definecolor{patchfigactive}{gray}{0.86}
\definecolor{patchfigoutline}{gray}{0.15}
\begin{tikzpicture}[x=0.63cm,y=0.66cm]
\path[use as bounding box] (-0.08,-0.15) rectangle (20.58,10.66);
\tikzset{
  patchfig site/.style={draw=patchfigoutline,line width=0.48pt},
  patchfig success/.style={draw=black,line width=0.86pt,
    -{Stealth[length=2.0mm,width=1.35mm]}},
  patchfig active/.style={fill=patchfigactive,draw=none},
  patchfig horizon/.style={draw=patchfigoutline,line width=0.72pt},
  patchfig boundary/.style={draw=patchfigoutline,line width=0.40pt},
  patchfig cut/.style={draw=patchfigoutline,line width=1.25pt,
    dash pattern=on 4.2pt off 2.6pt,line cap=round},
  patchfig small/.style={font=\scriptsize,text=patchfigoutline,inner sep=0pt}
}

% The striped vertical boundary used for every patch in Figure 1.  Drawing one
% continuous band along a site column is equivalent here because consecutive
% patches share the same vertical sides.
\newcommand{\ConvEndPatchBarrier}[3]{%
  \begin{scope}
    \clip (#1,#2) rectangle ({#1+1},#3);
    \foreach \k in {-1,...,68}{%
      \pgfmathsetmacro{\yy}{-0.045+0.165*\k}
      \fill[patchfigred]
        (#1,\yy)--(#1,{\yy+0.075})--({#1+0.105},{\yy+0.155})
        --({#1+0.105},{\yy+0.080})--cycle;
      \fill[patchfigred]
        ({#1+1},\yy)--({#1+1},{\yy+0.075})
        --({#1+0.895},{\yy+0.155})--({#1+0.895},{\yy+0.080})--cycle;}
  \end{scope}}

% A successful interaction cuts the source column and the two target columns.
\newcommand{\ConvEndInteractionBoundary}[2]{%
  \draw[patchfig boundary] ({#1-1},#2)--({#1+2},#2);}
\newcommand{\ConvEndInteractionMark}[2]{%
  \draw[patchfig success] ({#1+0.5},#2)--({#1-0.5},#2);
  \draw[patchfig success] ({#1+0.5},#2)--({#1+1.5},#2);
  \fill[black] ({#1+0.5},#2) circle[radius=1.25pt];}

% A patch ending at the source of a successful interaction has terminal type
% O.  As in Figure 1, these deterministically active patches are filled gray.
\foreach \x/\a/\b in {
  9/0.00/0.48,10/0.48/0.92,11/0.92/1.20,8/0.48/1.38,
  9/1.38/1.58,7/1.38/1.74,12/1.20/2.18,10/1.58/2.40,
  6/1.74/2.72,8/1.74/2.88,13/2.18/2.98,5/2.72/3.22,
  11/2.40/3.44,14/2.98/3.70,6/3.22/3.94,15/3.70/4.02,
  4/3.22/4.12,12/3.44/4.36,9/2.88/4.58,3/4.12/4.68,
  16/4.02/4.86,2/4.68/5.02,14/4.02/5.18,4/4.68/5.28,
  12/4.36/5.32,17/4.86/5.38,7/3.94/5.50,10/4.58/5.66,
  1/5.02/5.74,18/5.38/5.86,15/5.18/5.98,5/5.28/6.10
}{\path[patchfig active] (\x,\a) rectangle ({\x+1},\b);}

% Patch columns begin when their sites first enter the interaction cone.  The
% unequal starting times give the cone an irregular, asymmetric outline.
\foreach \x/\s in {
  0/5.74,1/5.02,2/4.68,3/4.12,4/3.22,
  5/2.72,6/1.74,7/1.38,8/0.48,9/0.00,
  10/0.00,11/0.92,12/1.20,13/2.18,14/2.98,
  15/3.70,16/4.02,17/4.86,18/5.38,19/5.86
}{\ConvEndPatchBarrier{\x}{\s}{10.50}}

% A site boundary begins when at least one of its neighboring site columns
% enters the cone.  No vertical grid is drawn where neither side is a patch.
\foreach \x/\s in {
  0/5.74,1/5.02,2/4.68,3/4.12,4/3.22,
  5/2.72,6/1.74,7/1.38,8/0.48,9/0.00,
  10/0.00,11/0.00,12/0.92,13/1.20,14/2.18,
  15/2.98,16/3.70,17/4.02,18/4.86,19/5.38,
  20/5.86
}{\draw[patchfig site](\x,\s)--(\x,10.50);}
\draw[patchfig horizon](0,0)--(20,0);
\draw[patchfig horizon](0,10.50)--(20,10.50);

% Dense successful-interaction skeleton below T.  The outer interactions grow
% the cone one site at a time; the remaining interactions subdivide its bulk.
\foreach \s/\y in {
  9/0.48,10/0.92,11/1.20,8/1.38,9/1.58,7/1.74,
  12/2.18,10/2.40,6/2.72,8/2.88,13/2.98,5/3.22,
  11/3.44,14/3.70,6/3.94,15/4.02,4/4.12,12/4.36,
  9/4.58,3/4.68,16/4.86,2/5.02,14/5.18,4/5.28,
  12/5.32,17/5.38,7/5.50,10/5.66,1/5.74,18/5.86,
  15/5.98,5/6.10
}{\ConvEndInteractionBoundary{\s}{\y}}

% Draw interactions last so that the black arrows and source dots remain
% legible over the gray fills and striped patch boundaries.
\foreach \s/\y in {
  9/0.48,10/0.92,11/1.20,8/1.38,9/1.58,7/1.74,
  12/2.18,10/2.40,6/2.72,8/2.88,13/2.98,5/3.22,
  11/3.44,14/3.70,6/3.94,15/4.02,4/4.12,12/4.36,
  9/4.58,3/4.68,16/4.86,2/5.02,14/5.18,4/5.28,
  12/5.32,17/5.38,7/5.50,10/5.66,1/5.74,18/5.86,
  15/5.98,5/6.10
}{\ConvEndInteractionMark{\s}{\y}}

% T cuts through every end patch: it is a guide, not a patch boundary.  There
% are no interactions or horizontal subdivisions between T and t.
\draw[patchfig cut] (0,6.35)--(20,6.35);
\node[patchfig small,anchor=west] at (20.18,0) {$0$};
\node[patchfig small,anchor=west] at (20.18,6.35) {$T$};
\node[patchfig small,anchor=west] at (20.18,10.50) {$t$};
\end{tikzpicture}
\endgroup
